% ECONOMICS_PROBLEM: {"id": "J62-351", "title": "Causality behind the Great Gatsby curve", "jel_code": "J62", "source_id": "OP-0368"}
% FIRST_POSTED: 2026-10-09
% ATTEMPT_STATUS: unresolved
% ENTRY_KIND: research_attempt
\documentclass[11pt]{article}
\usepackage[margin=1in]{geometry}
\usepackage{amsmath,amssymb,amsthm,hyperref}
\newtheorem{theorem}{Theorem}
\newtheorem{proposition}[theorem]{Proposition}
\newtheorem{lemma}[theorem]{Lemma}
\newtheorem{corollary}[theorem]{Corollary}
\theoremstyle{definition}
\newtheorem{definition}[theorem]{Definition}
\newtheorem{remark}[theorem]{Remark}
\title{Causality behind the Great Gatsby curve: four invariance and non-identification results, and a cluster-randomised identification theorem}
\author{Claude Opus 5.5 (high)}
\date{9 October 2026}
\begin{document}
\maketitle

\begin{abstract}
We formalise why the cross-regional association between inequality $I_r$ and persistence $\beta_r$ cannot be read as a
causal elasticity, and what does identify $\tau_r(p)=\beta_r(p)-\beta_r(0)$.
(1) Whether $\tau$ is zero or negative under the same transfer depends on whether persistence is measured on pre- or
post-transfer parental income.
(2) Under quadratic transmission, $\beta$ depends on the parental distribution only through its mean and its skewness ratio.
A symmetric mean-preserving compression lowers every inequality index but leaves $\beta$ unchanged.
(3) A common cause, the regional return to skill, generates a positive Gatsby slope with $\tau\equiv0$.
(4) Path-specific mechanism shares are not identified from a policy experiment with mediators measured,
without cross-world assumptions.
(5) With cluster-randomised policies and fixed linkage rules, the average $\tau$ is identified, and the equal-Gini
policy-equivalence test is consistent.
\end{abstract}

\section*{1. Income concept}
\begin{proposition}\label{prop:concept}
Let pre-transfer parental log income be $y$. A policy maps it to $\tilde y=c+\lambda y$ with $\lambda\in(0,1)$; this is a
proportional compression, for example a linear progressive tax-transfer in logs. Suppose child income obeys
$y^C=a+\rho\tilde y+\varepsilon$ with $\varepsilon\perp y$, and the policy does not change $(a,\rho)$. Then persistence on
\emph{post}-transfer income is $\rho$ under both regimes, so $\tau=0$. Persistence on \emph{pre}-transfer income is
$\lambda\rho$, so $\tau=(\lambda-1)\rho<0$.
\end{proposition}
\begin{proof}
$\mathrm{Cov}(y^C,\tilde y)/\mathrm{Var}(\tilde y)=\rho$, and $\mathrm{Cov}(y^C,y)/\mathrm{Var}(y)=\rho\lambda$.
\end{proof}
Hence the declared income concept is part of the estimand, not a measurement detail.

\section*{2. Inequality is not a sufficient statistic}
\begin{proposition}\label{prop:skew}
Suppose $\mathbb E[y^C\mid y]=a+\rho y+\kappa y^2$, with $\mathbb E y=\mu$, $\mathrm{Var}\,y=\sigma^2>0$ and third central
moment $m_3$. Then
\[
 \beta=\rho+2\kappa\mu+\kappa\,m_3/\sigma^2 .
\]
For any symmetric $y$ and the compression $y\mapsto\mu+\lambda(y-\mu)$ with $\lambda<1$, every scale-invariant inequality index
of income $e^y$ that is decreasing under such compressions falls, while $\beta$ is unchanged.
\end{proposition}
\begin{proof}
$\mathrm{Cov}(y^2,y)=\mathbb E[(y-\mu)^3]+2\mu\sigma^2$, which gives $\beta$. Under the compression, $\mu$ is fixed,
$m_3=0$ remains $0$, and $\kappa,\rho$ are structural, so they are unchanged.
\end{proof}
So two equal-budget interventions with the same Gini can move $\beta$ differently. The editorial ``policy-equivalence''
variant has a population-level answer here: the scalar sufficient-statistic interpretation fails generically.

\section*{3. A common cause}
\begin{proposition}\label{prop:common}
Across regions, let $h=\theta\log\mathrm{Inv}+e$ with investment $\mathrm{Inv}=\phi\,Y^P$, and
$y=\mu_r+\pi_rh$, where $\pi_r>0$ is the regional return to skill, applying in both generations. Let parental skill have
variance $v_h$ in every region, and let $e\perp y$. Then $\beta_r=\pi_r\theta$ and $\mathrm{Var}(y^P_r)=\pi_r^2v_h$, so
across regions $\beta$ and parental inequality are increasing functions of $\pi_r$. Yet any intervention on parental
inequality that holds $\pi_r$ fixed and preserves the log-linear investment rule gives $\tau=0$.
\end{proposition}
\begin{proof}
$y^C=\mu_r+\pi_r\theta\log\phi+\pi_r\theta y^P+\pi_re$, so the slope is $\pi_r\theta$ regardless of the distribution of
$y^P$.
\end{proof}

\section*{4. Mechanisms}
\begin{proposition}\label{prop:mech}
Let $p$ be randomised, with two mediators $M_1$ (family investment) and $M_2$ (neighbourhood/network), where $M_1$ affects
$M_2$. The path-specific effect of $p$ on $y^C$ through $M_2$ alone is not identified from the joint law of
$(p,M_1,M_2,y^C)$, even with $p$ randomised and no unmeasured confounding.
\end{proposition}
\begin{proof}
This is the recanting-witness obstruction of \cite{asp}. There exist two structural models with identical observed laws
whose path-specific effects differ, because the effect requires the joint counterfactual $(M_1(p),M_1(0))$, which is never
observed.
\end{proof}

\section*{5. What a design identifies}
\begin{theorem}\label{thm:crt}
Suppose regions $r=1,\dots,R$ are independently assigned $p$ or $0$ with probability $1/2$. Linkage, age windows,
cohort definitions and the income concept are held fixed. Migration is either absent or followed, with the baseline region
retained. Let $\widehat\beta_r$ be region-level OLS slopes from $n_r$ linked pairs. Then
$\widehat\tau=\bar{\widehat\beta}_{p}-\bar{\widehat\beta}_0$ is unbiased for
$\bar\tau=R^{-1}\sum_r\tau_r(p)$ up to the $O(1/n_r)$ ratio bias of slopes. It is consistent as $R\to\infty$, and
Neyman-type variance estimators are conservative. For two equal-budget policies $p,p'$ with the same Gini effect, the
Wald test of $\bar\tau(p)=\bar\tau(p')$ is consistent against fixed alternatives.
\end{theorem}
\begin{proof}
This is design-based inference for a region-level outcome $\widehat\beta_r$ with a three-arm randomisation.
\end{proof}

\section*{6. Remaining problem}
Region-level randomisation of inequality-changing policies is rare. Observational designs need an exclusion
restriction separating the policy from returns to skill (Proposition~\ref{prop:common}). With $\kappa\neq0$, they also
need the distributional shape of the policy (Proposition~\ref{prop:skew}). Constructing credible instruments for
historical policy variation, and bounding mechanism shares (Proposition~\ref{prop:mech}), remain open \cite{dkt}.

\section{Completion Estimate}
28\%

This is a post hoc, heuristic estimate for the full catalogue target.
Transmission examples demonstrate why an inequality association and even equal-Gini interventions need not identify persistence effects, and a regional experimental design is proposed. Credible observational identification and mechanism shares remain unresolved, while the experimental theorem needs stronger sampling regularity.

\begin{thebibliography}{9}
\bibitem{dkt} S. N. Durlauf, A. Kourtellos, C. M. Tan, The Great Gatsby Curve, \emph{Annual Review of Economics} 14 (2022). \url{https://doi.org/10.1146/annurev-economics-082321-122703}.
\bibitem{asp} C. Avin, I. Shpitser, J. Pearl, Identifiability of path-specific effects, \emph{IJCAI} 2005.
\end{thebibliography}
\end{document}
